\chapter{Triangulated Categories.}
\VerdierPhysicalPageAnchor{0110}%
\VerdierSemanticAnchor{ch2}
\begingroup
\small
\setlength{\abovedisplayskip}{4pt plus 2pt minus 2pt}%
\setlength{\belowdisplayskip}{4pt plus 2pt minus 2pt}%

\VerdierSemanticAnchor{ch2-1}
\section{Definitions and First Properties.}

\VerdierSemanticAnchor{ch2-1-1}
\subsection{Definition of Triangulated Categories.}

\VerdierSemanticAnchor{ch2-1-1-1}
\noindent\textbf{Definition 1.1.1.}\quad
A \emph{triangulated category} $\mathcal D$ is a strict additive
$\mathbb Z$-category (Chap. I, 1.5.11) endowed with a set of triangles
(Chap. I, 3.2.1), called \emph{distinguished triangles}, possessing the
following properties:

\VerdierSemanticAnchor{ch2-1-1-1-tri}
\smallskip
\noindent\textbf{TRI:}
Every triangle of $\mathcal D$ isomorphic to a distinguished triangle is a
distinguished triangle. For every object $X$ of $\mathcal D$, the triangle
\[
  X\xrightarrow{\operatorname{id}_X}X\longrightarrow 0
  \longrightarrow X[1]
\]
is distinguished. Every morphism $u\colon X\longrightarrow Y$ of
$\mathcal D$ is contained in a distinguished triangle
\[
  X\xrightarrow{u}Y\xrightarrow{v}Z\xrightarrow{w}X[1].
\]

\VerdierSemanticAnchor{ch2-1-1-1-trii}
\noindent\textbf{TRII:}
A triangle of $\mathcal D$
\[
  X\xrightarrow{u}Y\xrightarrow{v}Z\xrightarrow{w}X[1]
\]
is distinguished if and only if the triangle
\[
  Y\xrightarrow{v}Z\xrightarrow{w}X[1]
  \xrightarrow{-u[1]}Y[1]
\]
is distinguished.

\VerdierSemanticAnchor{ch2-1-1-1-triii}
\noindent\textbf{TRIII:}
For every pair of distinguished triangles

\VerdierSemanticAnchor{ch2-1-1-1-triii-triangles}
\[
  \begin{aligned}
    X  &\xrightarrow{u}  Y  \xrightarrow{v}  Z  \xrightarrow{w}  X[1],\\
    X' &\xrightarrow{u'} Y' \xrightarrow{v'} Z' \xrightarrow{w'} X'[1]
  \end{aligned}
\]
and every commutative diagram

\VerdierSemanticAnchor{ch2-1-1-1-triii-square}
\begin{center}
  \begin{tikzcd}[column sep=4.5em,row sep=3em]
    X \arrow[r,"u"] \arrow[d,"f"'] & Y \arrow[d,"g"] \\
    X' \arrow[r,"u'"] & Y'\,,
  \end{tikzcd}
\end{center}
\endgroup
